\begin{theorem}
Let \(s\to\infty\) and let \(a\) be a nonnegative integer such that \(a=o(s)\). Then
\[
  r(s,s+a)\ge (1+o(1))\frac{s}{e}\cdot
  2^{(s+a-1)/2-a^2/(2s)}.
\]
Equivalently, the formal statement uses the usual uniform \(\varepsilon\)-form for the \(1+o(1)\) lower factor when \(a/s\) is sufficiently small.
\end{theorem}
\begin{proof}
We prove the displayed \(1+o(1)\) assertion in the uniform
\(\varepsilon\)-form stated in the Lean theorem. Fix \(\varepsilon>0\). If
\(\varepsilon\ge1\), the factor \(1-\varepsilon\) is nonpositive, while
Ramsey numbers are nonnegative by \(\noderef{RamseyNumber}\); in this case
any positive choice of \(\delta\) and any \(s_0\ge1\) suffices. Hence assume
throughout that \(0<\varepsilon<1\).

Choose auxiliary constants
\[
  0<\rho<1,\qquad \eta>0,\qquad 0<\delta_1\le \frac12
\]
so small that
\[
  (1-\rho)2^{-\eta-\frac52\delta_1}\ge 1-\frac{\varepsilon}{2}. \tag{1}
\]
This is possible because the left hand side tends to \(1\) as
\(\rho,\eta,\delta_1\) tend to \(0\). Apply
\(\noderef{F2ForwardIndependentNearDiagonalBound}\) with the positive error
parameter \(\eta\). We obtain \(\delta_0>0\) and \(s_1\) such that, whenever
\(s\ge s_1\) and \(0\le a\le\delta_0s\), there is a loopless \(T_s\)-free
digraph with the exact F2 vertex count and the corresponding
forward-independent tuple bound. Set
\[
  \delta=\min(\delta_0,\delta_1).
\]
We will enlarge \(s_1\) to a final threshold \(s_0\). For \(s\ge4\), put
\[
  N_s=2^{2s-3}-2^{s-1}-2^{s-2}+1.
\]
Since
\[
  \frac{N_s}{2^{2s-3}}
  =1-2^{2-s}-2^{1-s}+2^{3-2s}\longrightarrow 1,
\]
we may require, by increasing \(s_0\), that
\[
  s_0\ge4,\qquad
  N_s\ge (1-\rho)2^{2s-3}\qquad(s\ge s_0). \tag{2}
\]
Also define
\[
  B_{s,a}=\frac{s}{e}\,
  2^{\frac{s+a-1}{2}-\frac{a^2}{2s}}.
\]
Because \(\delta\le1/2\), for \(0\le a\le\delta s\) we have
\[
  \frac{s+a-1}{2}-\frac{a^2}{2s}
  \ge \frac{s-1}{2}-\frac{\delta^2s}{2}
  \ge \frac{3s}{8}-\frac12,
\]
so \(B_{s,a}\to\infty\) uniformly in the range \(0\le a\le\delta s\).
Enlarge \(s_0\) again so that
\[
  B_{s,a}\ge \frac{2}{\varepsilon}
  \qquad(s\ge s_0,\;0\le a\le\delta s). \tag{3}
\]
Finally, since \(N_s\) grows like \(2^{2s}\) while
\[
  B_{s,a}\le \frac{s}{e}2^{(s+a-1)/2}
  \le \frac{s}{e}2^{3s/4}
\]
in the same range, enlarge \(s_0\) once more so that
\[
  N_s\ge (1-\varepsilon)B_{s,a}
  \qquad(s\ge s_0,\;0\le a\le\delta s). \tag{4}
\]

Now fix integers \(s,a\) with \(s\ge s_0\) and \(0\le a\le\delta s\), and
put
\[
  k=s+a.
\]
Since \(s_0\) was chosen at least \(4\), we have \(k\ge s\ge1\). Because
\(\delta\le\delta_0\), the hypotheses of
\(\noderef{F2ForwardIndependentNearDiagonalBound}\) are satisfied. Thus
there are a finite vertex set \(W\) and a loopless \(T_s\)-free digraph
\(D\) on \(W\) such that
\[
  |W|=N_s
\]
and
\[
  \overrightarrow{i_k}(D)\le
  2^E,\qquad
  E=\frac32s^2+as-\frac52s+\eta s. \tag{5}
\]
The looplessness and \(T_s\)-freeness are part of the same cited statement.

Apply \(\noderef{DigraphToGraphIndependentSetBound}\) to this \(D\), with
the above \(s\) and \(k\). It gives a \(K_s\)-free simple graph \(G\) on the
same vertex set \(W\) satisfying
\[
  i_k(G)\le \left(\frac e k\right)^k\overrightarrow{i_k}(D).
\]
Together with (5), this gives
\[
  i_k(G)\le \left(\frac e k\right)^k 2^E. \tag{6}
\]

First suppose \(i_k(G)=0\). By \(\noderef{SimpleGraphIndependentSetCount}\),
this means that \(G\) has no independent set of size \(k\). Together with
the \(K_s\)-freeness supplied by \(\noderef{DigraphToGraphIndependentSetBound}\),
the definition of \(\noderef{RamseyProperty}\) says that the
Ramsey property \(R(s,k,N_s)\) fails. If \(r(s,k)\le N_s\), then an induced
subgraph of \(G\) on \(r(s,k)\) vertices would also have neither a \(K_s\)
nor an independent \(k\)-set, contradicting the definition of
\(\noderef{RamseyNumber}\) as the least size at which
\(\noderef{RamseyProperty}\) holds. Hence \(r(s,k)>N_s\). By (4),
\[
  r(s,k)>(1-\varepsilon)B_{s,a},
\]
which is stronger than the required weak inequality.

It remains to handle the case \(i_k(G)>0\). Set
\[
  p=i_k(G)^{-1/k}.
\]
Since \(i_k(G)\) is a positive integer and \(k\ge1\), we have
\[
  0<p\le1,\qquad p^k i_k(G)=1.
\]
Therefore the hypotheses of \(\noderef{SamplingKsFreeRamseyBound}\) hold
for this \(G,s,k,p\), and it gives
\[
  pN_s-1<r(s,k). \tag{7}
\]
From (6) and positivity of \(i_k(G)\),
\[
  p=i_k(G)^{-1/k}\ge
  \frac{k}{e}\,2^{-E/k}. \tag{8}
\]
Combining (2) and (8), and using \(k=s+a\ge s\), we get
\[
\begin{aligned}
  pN_s
  &\ge
  (1-\rho)2^{2s-3}\frac{k}{e}2^{-E/k}  \\
  &\ge
  (1-\rho)\frac{s}{e}
  2^{\,2s-3-E/(s+a)}. \tag{9}
\end{aligned}
\]
We now expand the exponent in (9). Since \(k=s+a\),
\[
\begin{aligned}
  2s-3-\frac{E}{s+a}
  &=
  2s-3
  -\frac{\frac32s^2+as-\frac52s+\eta s}{s+a} \\
  &=
  \frac{s+a-1}{2}
  -\frac{a^2+5a}{2(s+a)}
  -\frac{\eta s}{s+a}. \tag{10}
\end{aligned}
\]
Comparing (10) with the target exponent gives
\[
\begin{aligned}
  &\left(\frac{s+a-1}{2}
  -\frac{a^2+5a}{2(s+a)}
  -\frac{\eta s}{s+a}\right)
  -
  \left(\frac{s+a-1}{2}-\frac{a^2}{2s}\right) \\
  &\qquad =
  \frac{a^3}{2s(s+a)}
  -\frac{5a}{2(s+a)}
  -\frac{\eta s}{s+a} \\
  &\qquad \ge
  -\frac{5}{2}\delta-\eta,
\end{aligned} \tag{11}
\]
because \(a\le\delta s\), \(s/(s+a)\le1\), and the first term is
nonnegative. Equations (9)--(11) imply
\[
  pN_s\ge
  (1-\rho)2^{-\frac52\delta-\eta}B_{s,a}.
\]
Since \(\delta\le\delta_1\), (1) gives
\[
  pN_s\ge \left(1-\frac{\varepsilon}{2}\right)B_{s,a}. \tag{12}
\]
Using (3), we may absorb the final \(-1\):
\[
  pN_s-1
  \ge
  \left(1-\frac{\varepsilon}{2}\right)B_{s,a}-1
  \ge
  (1-\varepsilon)B_{s,a}. \tag{13}
\]
Combining (7) and (13), and recalling \(k=s+a\), yields
\[
  (1-\varepsilon)\frac{s}{e}
  2^{\frac{s+a-1}{2}-\frac{a^2}{2s}}
  \le r(s,s+a).
\]
This is precisely the required uniform \(\varepsilon\)-form of the
close-to-diagonal bound.
\end{proof}
